\section{Claim Boundary Relative To The Proposal}

The proposal remains directionally correct, but the supported claim set is narrower than a full training-win story. This section makes the paper's claim boundary explicit.

\subsection{What Is Supported}

The current evidence supports the following:

\begin{itemize}
  \item trajectory-conditional necessity is a meaningful object on synthetic data;
  \item the same object survives larger-pool real-domain auditing in math;
  \item the project now has a reproducible object-validity protocol and a conservative Week~2 exit.
\end{itemize}

\subsection{What Is Not Yet Supported}

The current evidence does \emph{not} support:

\begin{itemize}
  \item a robust CNT-over-control offline training win;
  \item stable equal-budget utility gains;
  \item Week~4 ablations preserving a positive training advantage;
  \item equal-token frontier claims, second-domain transfer, or RL refinement.
\end{itemize}

\subsection{How Week 3 Should Now Be Read}

Stage~B should be read with a two-stage verdict:

\begin{enumerate}
  \item utility / no-harm first: CNT should not underperform matched control on original or damaged-prefix rollout utility;
  \item object-preservation second: weighted necessity and paraphrase/equiv stability can refine the reading once utility is non-negative.
\end{enumerate}

Under this rule, the current Stage~B families do not pass. The honest proposal-compatible interpretation is therefore:

\begin{quote}
object identified; audit protocol established; offline utility conversion unresolved.
\end{quote}

This is a clean negative / partial result, not a collapse of the full research direction.
